% \numberwithin{equation}{section}

\chapter{Conclusions}
Les pertorbacions de l'estat quiescent evolucionen segons els paràmetres del sistema com la constant d'acoblament,
la velocitat del so o la velocitat d'Alfvén, tant en l'interior de la protuberància com en l'exterior, en la corona.

La subdivisió del problema en les Eqs.~(\ref{eq.3.31}), (\ref{eq.3.32}) i (\ref{eq.3.33}) evoca dos tipus de solucions,
els modes d'Alfvén, molts semblants als d'una corda amb major densitat en el centre, i els modes ràpids i lents.
En el cas d'acoblament nul, $k_z = 0$, aquests darrers presenten la mateixa estructura que els modes d'Alfvén i, en conseqüència, que els de la corda.

Els autovalors i les autofuncions han estat calculades de manera correcta per Dedalus.
Aquest fet s'ha comprovat mitjançant el càlcul de l'error relatiu, $\eta$, respecte a les solucions analítiques amb forma de funcions trigonomètriques,
d'acord amb l'article de \textcite{oscillation_modes_II}.
L'error relatiu és molt menor a la unitat en els primers 100 autovalors, tot i que aquest creix significativament quan el nombre de l'autovalor augmenta. 
La situació observada també té lloc en el càlcul de les autosolucions d'una corda, i es deu a la manera de treballar de Dedalus.
Amb un nombre major d'autovalors calculats s'observa la mateixa tendència, conservant les proporcions.

L'anàlisi dels encreuaments evitats dels modes normals i l'acoblament d'aquests en modes ràpids i lents evidencia la importància de les condicions del sistema,
que modifiquen la constant d'acoblament i, en conseqüència, els modes normals que sorgeixen de fixar els extrems d'aquest.

L'estudi de les autofuncions és imprescindible a l'hora d'entendre el que succeeix en els encreuaments evitats dels modes normals,
és a dir, en els acoblaments dels modes ràpids i lents.
S'ha demostrat el canvi de dominància d'un mode a un altre en el sí d'un encreuament.
A més, s'ha estudiat l'evolució dels modes fonamentals híbrids per un ampli rang de valors de la constant d'acoblament.
D'aquesta manera es pot apreciar com els modes lents no es veuen molt afectats per un acoblament major,
mentre que els ràpids en pateixen les conseqüències de manera més notable, tant en el mode dominant com en el mode minoritari.

En el cas dels modes ràpids precisament, el mode minoritari, $\tilde{v}_{x}$, és aquell que presenta una direcció paral·lela al camp magnètic del sistema.
Per altra banda, el mode majoritari, $v_z$, és perpendicular al camp i paral·lel a l'acoblament.
Aquests sempre presenten paritats oposades, com s'ha demostrat en les figures.

Els modes lents en canvi, demostren un comportament oposat. Els seus modes minoritaris, $v_z$, són paral·lels a l'acoblament,
mentre que els predominants, $\tilde{v}_{x}$, són aquells perpendiculars al vector d'ona i paral·lels al camp.

La simulació de pertorbacions estacionàries en un estat quiescent d'una protuberància solar resulta en una descripció qualitativa del que té lloc en aquest fenòmen de l'atmosfera solar.
La sismologia de l'atmosfera solar és un objecte d'anàlisi que se'n pot beneficiar de projectes com el present.

Aquest és un estudi introductori d'acord amb les limitacions de temps i dificultat que suposa un treball de fi de grau.
Es pot ampliar aquest projecte amb l'anàlisi i desenvolupament dels períodes de les oscil·lacions, tema que s'ha descartat a causa de les limitacions mencionades.
Una altra manera d'augmentar la complexitat d'aquest treball consisteix en realitzar manco aproximacions al sistema de l'estudi,
com la d'ignorar els efectes centrífugs ò gravitatoris, entre d'altres.